\section{Related work and attribution}\label{sec:related-work}

\subsection{The shared tape model and the scope of the question}

This paper continues Brende and Claude Opus 5.5's \emph{Matrix Transposition
and the Tape Coding Conjecture}~\cite{FirstPaper2026}. We use its native
machine model, canonical entry width, cut inequality, and limiting path
condition. These are imported directly from the original
source package, together with its locality, counting and pass bounds and its
reduction to causal coding on tape networks.

Harvey and van der Hoeven's reduction relates multiplication to
transposition in the multitape model~\cite{HvdH2025}. Its coefficient-width
qualification matters here: their binary-matrix consequence of linear-time
multiplication has an iterated-logarithm loss, whereas their larger-coefficient
results yield the linear-time implication used for the canonical logarithmic
entry width. We do not replace the binary conclusion by a linear bound.
Classical transposition lower bounds for one-work-tape machines, including
the version with a separate output tape, concern restricted models
\cite{DMS1991,DietzfelbingerMaass1993}; they do not establish the unrestricted
fixed-multitape lower bound sought here.

Network coding has also been used to derive conditional multiplication
lower bounds~\cite{AFKL2019}. The general assertion that coding never improves
undirected multiple-unicast throughput has since been refuted
\cite{BravermanHe2026,ZhangLiLi2026}. Applying such a code to transposition would still require a tape
implementation with all preparation and access costs charged.

\subsection{Local coding, nonlinear recoding, and helper information}

Index coding provides the natural comparison for receivers that combine a
public message with their original rows. Bar-Yossef, Birk, Jayram, and Kol
relate index-code length to minrank and acyclic subgraphs~\cite{BBJK2011};
Blasiak, Kleinberg, and Lubetzky develop entropy-based bounds for broadcast
rate~\cite{BKL2011}. We also use established entropy methods, including Shearer's covering
inequality~\cite{CGFS1986}, and bounds obtained by counting code fibers.

Distributed and multi-sender index coding are closer comparisons because
sender knowledge constrains the permitted encoders. Sadeghi, Arbabjolfaei,
and Kim study distributed servers holding prescribed message subsets
\cite{SadeghiArbabjolfaeiKim2016}. Li, Ong, and Johnson develop cooperative
composite coding and a sender-constrained rank-minimization formulation
\cite{LiOngJohnson2019,LiOngJohnsonRank2019}. Our fixed original-row packet
models can be interpreted in this setting after specifying the entry-level
messages and column demands. Adaptive queries additionally choose their
access supports from the public history; a fixed sender allocation does not
by itself express this interface.

The adaptive separator theorem addresses the following interface: each query
has specified access to \emph{original} rows, its support and subsequent
queries may depend on the public history, and the correct-input law can be
correlated and nonuniform. Its graph records which rows a query accesses
jointly; it is not the receiver side-information graph of standard index
coding. The theorem charges both the helper and source entropy deficit and
allows a separator chosen from the completed public history. The original-row access and public-history hypotheses distinguish this
statement from the cited index-coding bounds.

Even multi-sender results require care. In the single-uniprior model of
Ong, Ho, and Lim~\cite{OngHoLim2016}, each receiver initially knows one
message, known to no other receiver, and its demands consist of whole
messages. Bundling each row into a message changes our column demands into
partial-message demands; splitting the rows into entries changes the
single-uniprior hypothesis. Thus that characterization cannot simply be
substituted for the original-row locality argument.

Our second family permits arbitrary bijective nonlinear recoding within
rows but restricts the subsequent packets to fixed single-coordinate or
pairwise XORs. This restriction is substantial: nonlinear index codes can
outperform linear ones~\cite{LubetzkyStav2008}. The component-flip and graph-distance arguments apply to this fixed-packet
architecture. Similarly, the two-stage helper theorem counts fibers of a specific
triangular--causal composition. Its helper selects the whole composition; no additional uncharged
global operation is allowed between the stages.

\subsection{Affine geometry, expanders, and paid computation}

There is a longstanding distinction between simple one-dimensional affine
constructions and expansion. Klawe ruled out several proposed expander
constructions, including one-dimensional analogues of the Gabber--Galil
construction~\cite{Klawe1984}. Our affine-motion theorem has a different
statement: it concerns countably many local affine branches of normalized
head paths, measures exceptional time, and applies the original limiting
cut condition. The approximation version retains the stated subsequence
and bounded-cover hypotheses. We do not identify these conclusions with
Klawe's finite-graph theorem.

Monotone expanders and their connections to page and pushdown graphs were
studied by Dvir and Wigderson~\cite{DvirWigderson2010}.
Bourgain and Yehudayoff constructed a continuous monotone expander using
M\"obius transformations and proved the associated growth and restricted
spectral-gap results~\cite{BourgainYehudayoff2013}. Our analytic development follows this established work, proving the
rational-family, growth, flattening and spectral statements needed for the
native implementation.

The scalar sum-product and projection ingredients belong to Bourgain's
discretized ring and projection theory~\cite{Bourgain2003,Bourgain2010}.
Their mathematical origin is distinct from the formal source we reuse.
The pinned OpenAI library~\cite{OAIProjection2026} supplies more than the
endpoint \lean{scalar_weak_radius_projection}: we directly import
\lean{ScalarDilationBound}, \lean{ScalarPluennecke},
\lean{ScalarBallCandidates}, \lean{ScalarCoverCount}, and
\lean{ScalarWeakRadiusProjection}. The bridge also uses its scalar-cell,
ball-mass, zoom, and ratio-family definitions and estimates, including
\lean{scalar_zoom_cells}, \lean{scalar_zoom_ball_bound},
\lean{scalar_refined_ball_bound},
\lean{scalar_ratio_family_ball}, and
\lean{scalar_ratio_family_card}.
The local arguments verify the hypotheses and derive the expansion inputs.
The bibliography pins the formal source revision, whose imported proofs are
checked in the same dependency chain.

The padding theorem adds a computational conclusion: the walk is generated,
executed, cleaned up, and followed by the original computation at additive
linear cost, with the stated fixed-machine and resource quantifiers.
These costs do not follow merely from a graph's expansion or a formula for
its maps. The scan-palette and recursion exclusions also account for actual
execution, under the stated affine regularity and exceptional-time hypotheses.

\subsection{The transpose prediction conjecture is already false}

Conjecture~36 of Mittal and Raz~\cite{MittalRaz2021} proposed an exponentially
small success bound for a matrix-transpose prediction game. Its failure is
explicitly recorded in Mittal's dissertation~\cite[\S10.8.2, Proposition~10.8.3
and Lemma~10.8.5]{Mittal2025}, which credits a proof of Braverman and Forbes
from their 2022 personal communication. We cite the dissertation as the
accessible proof and claim no new refutation.

Related counterexamples also follow from earlier coded-caching constructions:
placement delivery arrays of Yan et al.~\cite{YanEtAl2017}, induced-matching
schemes of Shanmugam, Tulino, and Dimakis~\cite{ShanmugamEtAl2017}, and the
projective-geometric schemes of Chittoor et al.~\cite{ChittoorEtAl2021}.
These papers do not themselves state a disproof of Conjecture~36.
Their induced-matching classes give XOR prediction strategies with success
probability $2^{-r}$, where $r$ is the number of nonempty classes. Suitable
square conversions give $o(d)$ observed rows per player and $r=o(d^2)$.
For the dense-graph route, the construction of Alon, Moitra, and
Sudakov~\cite{AMS2013} is combined with minimum-degree pruning. This avoids a
fixed-vertex degree estimate in the supplied v1 caching preprint. The
conversions and version-specific caveat are detailed in
\lean{paper/prior-art.txt}; they are background arguments, not new Lean
results claimed here.

The game does not charge for acquiring or combining the observed rows.
Its counterexamples close that prediction-game route to a lower bound but
do not supply a causal tape implementation.

\subsection{Formal dependencies and novelty boundaries}

The formal proofs use Lean~4~\cite{Lean4} and the pinned Mathlib4 library
\cite{Mathlib2020,Mathlib4}, in addition to the two explicit source dependencies
above. Standard finite-set, entropy, measure, topology, and arithmetic
results are reused. We distinguish proving the present endpoints from
claiming new underlying analytic or information-theoretic principles.

The literature comparison covers the cited primary sources and targeted
searches for local index coding, two-stage causal maps, affine motion and
recursive tape schedules. It does not establish exhaustive priority. We
claim the stated local-coding bounds, motion consequences and native
implementations, while crediting their analytic and information-theoretic
ingredients. Supporting lemmas are included for their role in these arguments,
without a blanket novelty claim.
